\documentclass[../main.tex]{subfiles} % 使用 subfiles 文档类，指定主文档
\graphicspath{{\subfix{../image/}}} % 指定图片目录，后续可以直接使用图片文件名。

% 例如:
% \begin{figure}[H]
% \centering
% \includegraphics[width=0.3\textwidth]{图.png}
% \caption{}
% \label{figure:图}
% \end{figure}
% 注意:上述\label{}一定要放在\caption{}之后,否则引用图片序号会只会显示??.

\begin{document}

\chapter{热方程}

\begin{definition}
设 $\Omega\subset\mathbb{R}^n$，$a>0$ 为常数，$f=f(x,t)$ 为已知函数，$u=u(x,t)$，$x\in\Omega$，$t>0$. 称方程
\begin{align}
u_t - a^2\Delta u = f \label{def:heat}
\end{align}
为\textbf{热方程}. 当 $f\equiv 0$ 时，称
\begin{align}
u_t - a^2\Delta u = 0 \label{def:heat-hom}
\end{align}
为\textbf{齐次热方程}.
\end{definition}

\begin{definition}
设 $Q\subset\mathbb{R}^n$，$T>0$. 我们用$C^{2,1}(Q)$表示所有在$Q$内关于$x$二阶偏导数连续,关于$t$一阶偏导数连续的函数构成的函数集,即
\begin{align*}
C^{2,1}(Q) \triangleq \{\, u \mid u, u_t, u_{x_i}, u_{x_ix_j} \in C(Q),\ i,j=1,\cdots,n\,\}.
\end{align*}
称热方程 \eqref{def:heat}在$C^{2,1}$函数集中的解为热方程\eqref{def:heat}的\textbf{古典解}. 类似地，用$C^{1,0}(Q)$表示所有在$Q$内函数本身连续且关于$x$一阶偏导数连续的函数构成的函数集,即
\begin{align*}
C^{1,0}(Q) \triangleq \{\, u \mid u, u_{x_i} \in C(Q),\ i=1,\cdots,n\,\}.
\end{align*}
\end{definition}

\begin{definition}
热方程 \eqref{def:heat} 的定解条件包括初始条件和边值条件：
\begin{enumerate}[(1)]
\item \textbf{初始条件：}
\begin{align*}
u(x,0)=\varphi(x),\qquad x\in\Omega,
\end{align*}
其中 $\varphi$ 为已知函数.称\eqref{eq:heat-cauchy1d}式
\begin{align}
\begin{cases}
\frac{\partial u}{\partial t} - a^2\frac{\partial^2 u}{\partial x^2} = f(x,t), & (x,t)\in\mathbb{R}\times\mathbb{R}_+, \\
u(x,0)=\varphi(x), & x\in\mathbb{R}.
\end{cases} \label{eq:heat-cauchy1d}
\end{align}
为\textbf{一维热方程初值问题}.称\eqref{eq:多维热方程初值问题}式
\begin{align}\label{eq:多维热方程初值问题}
\begin{cases}
\frac{\partial u}{\partial t}-a^2\Delta u=f(x,t), & (x,t)\in\mathbb{R}_+^{n+1}=\mathbb{R}^n\times\mathbb{R}_+,\\
u(x,0)=\varphi(x), & x\in\mathbb{R}^n.
\end{cases}
\end{align}
为\textbf{$n$维热方程初值问题}.

\item \textbf{边值条件：} 通常有以下三类：
\begin{enumerate}[(i)]
\item \textbf{第一类边值条件（Dirichlet 条件）：}
\begin{align*}
u(x,t)=g(x,t),\qquad x\in\partial\Omega,\ t>0.
\end{align*}
当 $g\equiv$ 常数时，称边界保持恒温.

\item \textbf{第二类边值条件（Neumann 条件）：}
\begin{align*}
k\frac{\partial u}{\partial n}(x,t)=g(x,t),\qquad x\in\partial\Omega,\ t>0,
\end{align*}
其中 $n$ 为 $\partial\Omega$ 的单位外法向量，$k>0$ 为导热系数. 当 $g\equiv 0$ 时，称边界绝热.

\item \textbf{第三类边值条件（Robin 条件）：}
\begin{align*}
k\frac{\partial u}{\partial n}(x,t)=\alpha_0(x,t)(g_0(x,t)-u(x,t)),\qquad x\in\partial\Omega,\ t>0,
\end{align*}
或等价地
\begin{align*}
\frac{\partial u}{\partial n}(x,t)+\alpha(x,t)u(x,t)=g(x,t),\qquad x\in\partial\Omega,\ t>0,
\end{align*}
其中 $\alpha_0>0$ 为热交换系数，$g_0$ 为周围介质温度，$\alpha=\frac{\alpha_0}{k}$，$g=\frac{\alpha_0 g_0}{k}$.
\end{enumerate}
\end{enumerate}
\end{definition}
\begin{remark}
热方程 \eqref{def:heat} 的物理背景：若 $u(x,t)$ 表示物体在位置 $x$、时刻 $t$ 的温度，则 $a=\sqrt{k/(c\rho)}$，其中 $k$ 为导热系数，$c$ 为比热，$\rho$ 为密度，$f=f_0/c$，$f_0$ 为热源密度. 若 $u$ 表示浓度，则 \eqref{def:heat} 称为反应扩散方程，此时 $a^2$ 为扩散系数，$-a^2\Delta u$ 为扩散项，$f$ 为反应项.
\end{remark}

\subfile{Chapter-03/section-01.tex}

\subfile{Chapter-03/section-02.tex}

\subfile{Chapter-03/section-03.tex}

\subfile{Chapter-03/section-04.tex}

\subfile{Chapter-03/section-05.tex}

\subfile{Chapter-03/section-06.tex}

\subfile{Chapter-03/section-07.tex}

\subfile{Chapter-03/section-08.tex}

\subfile{Chapter-03/section-09.tex}

\subfile{Chapter-03/section-10.tex}

\end{document}